\documentclass[../Main.tex]{subfiles}

\begin{document}
\section{Introduction}
We will consider flows of the form:
\begin{equation*}
    \vec{u} = u(y, t) \vec{e_x}
\end{equation*}
called 2D parallel viscous flow.

\begin{definition}{Viscous stress}
    The \underline{viscous stress} is a shear stress given by:
    \begin{equation*}
        \tau = \mu \frac{\partial u}{\partial \vec{n}}
    \end{equation*}
    where $\mu$ is a property of the fluid called the (dynamic) viscosity.
\end{definition}
In practice, $\vec{n} = \pm \vec{e_y}$.
\section{Forces}
The force balance along $x$ and $y$ are:
\begin{align*}
    \rho \frac{\partial u}{\partial t} &= \mu \frac{\partial^{2}u}{\partial y^{2}} - \frac{\partial p}{\partial x} + f_x \\
    0 &= -\frac{\partial p}{\partial y} + f_y
\end{align*}
so in the $y$ direction we just have a hydrostatic balance.

The boundary condition in the case of viscous flows is $\vec{u} = \vec{U}$ (no-slip condition), using the notation from earlier.

At an interface between two fluids of different viscosities, we must match the stresses:
\begin{equation*}
    \mu_1 \frac{\partial u_1}{\partial y} = \mu_2 \frac{\partial u_2}{\partial y}.
\end{equation*}
\section{Unsteady Flow}
To solve problems involving unsteady flow, we introduce the \underline{kinematic viscosity} $\nu = \frac{\mu}{\rho}$. Then we find that the equation for $u$ reduces to a diffusion equation:
\begin{equation*}
    \frac{\partial u}{\partial t} = \nu \frac{\partial^{2}u}{\partial y^{2}}.
\end{equation*}
And we can solve this using either dimensional analysis or a Fourier transform.
\end{document}